% GENERATED FILE. Edit canonical lessons, then run npm run generate:books.
% canonical-source-hash: fnv1a64:2e58051bb315c348
% canonical-lessons: TA-C305-nirupar, TA-C305-neralai, TA-C305-oliparappu, TA-C305-tokuppalar, TA-C305-olippativalar

\chapter{On Air}
\label{ch:ta-on-air}
\hlchaptermodality{305}

\begin{chapteropening}
\textbf{By the end of this chapter:} \emph{I can talk about news on television and radio.}

Complete the last lesson of chapter 305: I can talk about news on television and radio.
\end{chapteropening}

\section[\texorpdfstring{\ta{நிருபர்} (nirupar)}{நிருபர் (nirupar)}]{\ta{நிருபர்} (nirupar) — a reporter}

\label{lesson:TA-C305-nirupar}



\begin{quote}
Before the new one: say the Tamil for to shake up, then the Tamil for to ruin.
\end{quote}

\subsection*{You'll want to know: \ta{நிருபர்}}
\textbf{\ta{நிருபர்}} — \emph{nirupar} — \textquotedblleft{}a reporter\textquotedblright{}.

\begin{quote}
\ta{ஒரு} \ta{நிருபர்} \ta{விபத்து} \ta{நடந்த} \ta{இடத்துக்கு} \ta{வந்தார்}.
\end{quote}

\begin{quote}
\emph{oru nirupar vipattu naṭanta iṭattukku vantār.}
\end{quote}

\begin{quote}
A reporter came to the place where the accident happened.
\end{quote}

Plural \ta{நிருபர்கள்}.

\begin{cousinweb}
\textbf{\ta{நிருபர்}} (\emph{nirupar}), \textquotedblleft{}a reporter, a correspondent\textquotedblright{}, is a Tamil word whose origin is not settled in the sources at hand.
\end{cousinweb}

\begin{grammarlens}[title={What you've built}]
One more word for telling what happened.
\end{grammarlens}

\subsection*{Your turn}
\begin{itemize}
\raggedright
  \item \emph{Say it:} \emph{nirupar}
  \item \emph{Say it:} \emph{nirupar}, once more
  \item \emph{Say it:} \emph{oru nirupar vipattu naṭanta iṭattukku vantār.} — A reporter came to the place where the accident happened.
  \item \emph{Recall:} say the Tamil for to shake up, then the Tamil for to ruin, then say \emph{nirupar} again
\end{itemize}

\begin{tcolorbox}[breakable,colback=teal!4,colframe=teal!35!black,title={Before you move on}]
What does \ta{நிருபர்} mean? (\textquotedblleft{}A reporter\textquotedblright{}.) Say it once more.
\end{tcolorbox}
\section[\texorpdfstring{\ta{நேரலை} (nēralai)}{நேரலை (nēralai)}]{\ta{நேரலை} (nēralai) — a live broadcast}

\label{lesson:TA-C305-neralai}



\begin{quote}
Before the new one: say the Tamil for to ruin, then the Tamil for a reporter.
\end{quote}

\subsection*{You'll want to know: \ta{நேரலை}}
\textbf{\ta{நேரலை}} — \emph{nēralai} — \textquotedblleft{}a live broadcast\textquotedblright{}.

\begin{quote}
\ta{தொலைக்காட்சியில்} \ta{போட்டியை} \ta{நேரலையில்} \ta{பார்த்தோம்}.
\end{quote}

\begin{quote}
\emph{tolaikkāṭciyil pōṭṭiyai nēralaiyil pārttōm.}
\end{quote}

\begin{quote}
We watched the match live on television.
\end{quote}

\ta{நேரலையில்}, live (on air).

\begin{cousinweb}
\textbf{\ta{நேரலை}} (\emph{nēralai}) is a Tamil compound built from \ta{நேர்}, \textquotedblleft{}straightness, directness\textquotedblright{} (DEDR 3772), and \ta{அலை}, \textquotedblleft{}a wave\textquotedblright{}.
\end{cousinweb}

\begin{grammarlens}[title={What you've built}]
One more word for telling what happened.
\end{grammarlens}

\subsection*{Your turn}
\begin{itemize}
\raggedright
  \item \emph{Say it:} \emph{nēralai}
  \item \emph{Say it:} \emph{nēralai}, once more
  \item \emph{Say it:} \emph{tolaikkāṭciyil pōṭṭiyai nēralaiyil pārttōm.} — We watched the match live on television.
  \item \emph{Recall:} say the Tamil for to ruin, then the Tamil for a reporter, then say \emph{nēralai} again
\end{itemize}

\begin{tcolorbox}[breakable,colback=teal!4,colframe=teal!35!black,title={Before you move on}]
What does \ta{நேரலை} mean? (\textquotedblleft{}A live broadcast\textquotedblright{}.) Say it once more.
\end{tcolorbox}
\section[\texorpdfstring{\ta{ஒளிபரப்பு} (oḷiparappu)}{ஒளிபரப்பு (oḷiparappu)}]{\ta{ஒளிபரப்பு} (oḷiparappu) — to broadcast}

\label{lesson:TA-C305-oliparappu}



\begin{quote}
Before the new one: say the Tamil for a reporter, then the Tamil for a live broadcast.
\end{quote}

\subsection*{You'll want to know: \ta{ஒளிபரப்பு}}
\textbf{\ta{ஒளிபரப்பு}} — \emph{oḷiparappu} — \textquotedblleft{}to broadcast\textquotedblright{}.

\begin{quote}
\ta{தொலைக்காட்சி} \ta{அந்தச்} \ta{செய்தியை} \ta{உடனே} \ta{ஒளிபரப்பியது}.
\end{quote}

\begin{quote}
\emph{tolaikkāṭci antac ceytiyai uṭaṉē oḷiparappiyatu.}
\end{quote}

\begin{quote}
Television broadcast that news at once.
\end{quote}

Past (I …): \ta{ஒளிபரப்பினேன்}, \emph{oḷiparappiṉēṉ}.

\ta{ஒளி}, light + \ta{பரப்பு}, to spread: said of television.

\begin{cousinweb}
\textbf{\ta{ஒளிபரப்பு}} (\emph{oḷiparappu}) is a Tamil compound built from \ta{ஒளி}, \textquotedblleft{}light, brightness\textquotedblright{} (DEDR 1016), and \ta{பரப்பு}, \textquotedblleft{}to spread, disseminate news\textquotedblright{} (DEDR 3949).
\end{cousinweb}

\begin{grammarlens}[title={What you've built}]
One more word for telling what happened.
\end{grammarlens}

\subsection*{Your turn}
\begin{itemize}
\raggedright
  \item \emph{Say it:} \emph{oḷiparappu}
  \item \emph{Say it:} \emph{oḷiparappu}, once more
  \item \emph{Say it:} \emph{tolaikkāṭci antac ceytiyai uṭaṉē oḷiparappiyatu.} — Television broadcast that news at once.
  \item \emph{Recall:} say the Tamil for a reporter, then the Tamil for a live broadcast, then say \emph{oḷiparappu} again
\end{itemize}

\begin{tcolorbox}[breakable,colback=teal!4,colframe=teal!35!black,title={Before you move on}]
What does \ta{ஒளிபரப்பு} mean? (\textquotedblleft{}To broadcast\textquotedblright{}.) Say it once more.
\end{tcolorbox}
\section[\texorpdfstring{\ta{தொகுப்பாளர்} (tokuppāḷar)}{தொகுப்பாளர் (tokuppāḷar)}]{\ta{தொகுப்பாளர்} (tokuppāḷar) — a presenter; a host}

\label{lesson:TA-C305-tokuppalar}



\begin{quote}
Before the new one: say the Tamil for a live broadcast, then the Tamil for to broadcast.
\end{quote}

\subsection*{You'll want to know: \ta{தொகுப்பாளர்}}
\textbf{\ta{தொகுப்பாளர்}} — \emph{tokuppāḷar} — \textquotedblleft{}a presenter; a host\textquotedblright{}.

\begin{quote}
\ta{தொகுப்பாளர்} \ta{நிருபரிடம்} \ta{ஒரு} \ta{கேள்வி} \ta{கேட்டார்}.
\end{quote}

\begin{quote}
\emph{tokuppāḷar nirupariṭam oru kēḷvi kēṭṭār.}
\end{quote}

\begin{quote}
The presenter asked the reporter a question.
\end{quote}

-\ta{ஆளர்} marks one who does a job: the one who puts a programme together.

\begin{cousinweb}
\textbf{\ta{தொகுப்பாளர்}} (\emph{tokuppāḷar}) is built from \ta{தொகுப்பு}, \textquotedblleft{}a putting together\textquotedblright{}, from native \ta{தொகு}, \textquotedblleft{}to collect, assemble\textquotedblright{} (DEDR 3476), with -\ta{ஆளர்}, \textquotedblleft{}one who\textquotedblright{}.
\end{cousinweb}

\begin{grammarlens}[title={What you've built}]
One more word for telling what happened.
\end{grammarlens}

\subsection*{Your turn}
\begin{itemize}
\raggedright
  \item \emph{Say it:} \emph{tokuppāḷar}
  \item \emph{Say it:} \emph{tokuppāḷar}, once more
  \item \emph{Say it:} \emph{tokuppāḷar nirupariṭam oru kēḷvi kēṭṭār.} — The presenter asked the reporter a question.
  \item \emph{Recall:} say the Tamil for a live broadcast, then the Tamil for to broadcast, then say \emph{tokuppāḷar} again
\end{itemize}

\begin{tcolorbox}[breakable,colback=teal!4,colframe=teal!35!black,title={Before you move on}]
What does \ta{தொகுப்பாளர்} mean? (\textquotedblleft{}A presenter; a host\textquotedblright{}.) Say it once more.
\end{tcolorbox}
\section[\texorpdfstring{\ta{ஒளிப்பதிவாளர்} (oḷippativāḷar)}{ஒளிப்பதிவாளர் (oḷippativāḷar)}]{\ta{ஒளிப்பதிவாளர்} (oḷippativāḷar) — a cameraman}

\label{lesson:TA-C305-olippativalar}



\begin{quote}
Before the new one: say the Tamil for to broadcast, then the Tamil for a presenter.
\end{quote}

\subsection*{You'll want to know: \ta{ஒளிப்பதிவாளர்}}
\textbf{\ta{ஒளிப்பதிவாளர்}} — \emph{oḷippativāḷar} — \textquotedblleft{}a cameraman\textquotedblright{}.

\begin{quote}
\ta{நிருபருடன்} \ta{ஒரு} \ta{ஒளிப்பதிவாளரும்} \ta{வந்தார்}.
\end{quote}

\begin{quote}
\emph{niruparuṭaṉ oru oḷippativāḷarum vantār.}
\end{quote}

\begin{quote}
A cameraman came along with the reporter.
\end{quote}

\ta{ஒளிப்பதிவு}, video recording + -\ta{ஆளர்}, one who.

\begin{cousinweb}
\textbf{\ta{ஒளிப்பதிவாளர்}} (\emph{oḷippativāḷar}) is built from \ta{ஒளி}, \textquotedblleft{}light\textquotedblright{} (DEDR 1016), and \ta{பதிவு}, \textquotedblleft{}an impression, a record\textquotedblright{} (DEDR 3911), with -\ta{ஆளர்}, \textquotedblleft{}one who\textquotedblright{}.
\end{cousinweb}

\begin{grammarlens}[title={What you've built}]
One more word for telling what happened.
\end{grammarlens}

\subsection*{Your turn}
\begin{itemize}
\raggedright
  \item \emph{Say it:} \emph{oḷippativāḷar}
  \item \emph{Say it:} \emph{oḷippativāḷar}, once more
  \item \emph{Say it:} \emph{niruparuṭaṉ oru oḷippativāḷarum vantār.} — A cameraman came along with the reporter.
  \item \emph{Recall:} say the Tamil for to broadcast, then the Tamil for a presenter, then say \emph{oḷippativāḷar} again
\end{itemize}

\begin{tcolorbox}[breakable,colback=teal!4,colframe=teal!35!black,title={Before you move on}]
What does \ta{ஒளிப்பதிவாளர்} mean? (\textquotedblleft{}A cameraman\textquotedblright{}.) Say it once more.
\end{tcolorbox}
